\documentclass[10pt,aspectratio=169]{beamer}
\input{../../_en_preambulo_beamer}% Comandos y macros estadísticas compartidas para las presentaciones Beamer en inglés (Metropolis)

% Conjuntos numéricos básicos
\newcommand{\R}{\mathbb{R}}
\newcommand{\N}{\mathbb{N}}
\newcommand{\Q}{\mathbb{Q}}
\newcommand{\Z}{\mathbb{Z}}
\newcommand{\set}[1]{\left\{ #1 \right\}}
\newcommand{\abs}[1]{\left|#1\right|}
\newcommand{\norm}[1]{\left\| #1 \right\|}

% Comandos estadísticos y probabilísticos del curso
\newcommand{\Var}{\operatorname{Var}}
\newcommand{\cov}{\operatorname{Cov}}
\newcommand{\corr}{\rho}
\newcommand{\comb}[2]{\begin{pmatrix} #1 \\ #2 \end{pmatrix}}
\newcommand{\s}{\sigma}
\newcommand{\particion}{\mathfrak{P}}
\newcommand{\card}[1]{\#\left( #1 \right)}
\newcommand{\seq}[3]{\set{#1}_{#2}^{#3}}
\newcommand{\E}{\mathbb{E}}
\newcommand{\Prob}{\mathbb{P}}

% Entornos pedagógicos y cajas en inglés académico natural (soportando tanto nombres en inglés como en español)
\newenvironment{discussion}{%
  \begin{alertblock}{Discussion Question}%
}{%
  \end{alertblock}%
}
\newenvironment{discusion}{%
  \begin{alertblock}{Discussion Question}%
}{%
  \end{alertblock}%
}

\newenvironment{reflection}{%
  \begin{alertblock}{Classroom Reflection}%
}{%
  \end{alertblock}%
}
\newenvironment{reflexion}{%
  \begin{alertblock}{Classroom Reflection}%
}{%
  \end{alertblock}%
}

\newenvironment{paradox}{%
  \begin{alertblock}{Exploratory Paradox}%
}{%
  \end{alertblock}%
}
\newenvironment{paradoja}{%
  \begin{alertblock}{Exploratory Paradox}%
}{%
  \end{alertblock}%
}

\newenvironment{motivation}{%
  \begin{exampleblock}{Motivation: Data Science in Practice}%
}{%
  \end{exampleblock}%
}
\newenvironment{motivacion}{%
  \begin{exampleblock}{Motivation: Data Science in Practice}%
}{%
  \end{exampleblock}%
}

\newenvironment{quickchallenge}{%
  \begin{block}{Quick Team Challenge}%
}{%
  \end{block}%
}
\newenvironment{retoclase}{%
  \begin{block}{Quick Team Challenge}%
}{%
  \end{block}%
}

\title[Fixed-effect effects]{Section 08.03: Fixed-Effect Model Effects}\subtitle{Mean squares and the $F$ statistic}
\author[J. Castillo Colmenares]{Juliho Castillo Colmenares}\institute{Tecnológico de Monterrey}\date{}
\begin{document}
\begin{frame}[plain]\titlepage\end{frame}
\begin{frame}{Roadmap}\small\begin{enumerate}\item Mean squares.\item Hypotheses and $F$ statistic.\item ANOVA table and interpretation.\end{enumerate}\end{frame}
\begin{frame}{Source and scope}\small This section explains how treatment effects change the mean squares in a one-factor ANOVA.\begin{block}{Guiding question}Why can a mean-square ratio detect effects?\end{block}\end{frame}
\begin{frame}{Mean squares}\small\[\text{CMTR}=\frac{\text{SCTR}}{k-1},\qquad \text{CME}=\frac{\text{SCE}}{N-k}.\]Each sum of squares is divided by its degrees of freedom.\end{frame}
\begin{frame}{Expected mean squares}\small\[E(\text{CME})=\sigma^2,\qquad E(\text{CMTR})=\sigma^2+\frac{\sum_i n_i\tau_i^2}{k-1}.\]\end{frame}
\begin{frame}{Omnibus hypotheses}\small\[H_0:\tau_1=\cdots=\tau_k=0\quad\Longleftrightarrow\quad\mu_1=\cdots=\mu_k.\]\[H_a:\text{at least one effect is nonzero}.\]\end{frame}
\begin{frame}{Ratio logic}\small Under $H_0$, CMTR and CME estimate the same variance. Under $H_a$, CMTR contains a positive treatment term.\end{frame}
\begin{frame}{$F$ statistic}\small\[F=\frac{\text{CMTR}}{\text{CME}}=\frac{\text{SCTR}/(k-1)}{\text{SCE}/(N-k)}.\]\end{frame}
\begin{frame}{Null distribution}\small Under the model assumptions and $H_0$,\[F\sim F_{k-1,N-k}.\]It is central when all effects vanish.\end{frame}
\begin{frame}{Decision rule}\small Reject $H_0$ at level $\alpha$ when\[F>F_{\alpha,k-1,N-k}\]or, equivalently, when the $p$-value is below $\alpha$.\end{frame}
\begin{frame}{ANOVA table}\small\begin{itemize}\item Treatments: SCTR, $k-1$, CMTR.\item Error: SCE, $N-k$, CME.\item Total: SCT, $N-1$.\end{itemize}\end{frame}
\begin{frame}{Procedure I}\small\begin{enumerate}\item Check randomization, independence, and approximate error normality.\item Compute means and sums of squares.\item Obtain CMTR and CME.\end{enumerate}\end{frame}
\begin{frame}{Procedure II}\small\begin{enumerate}\setcounter{enumi}{3}\item Form $F$.\item Compare with a quantile or compute $p$.\item Report the global conclusion.\end{enumerate}\end{frame}
\begin{frame}{Application: treatments}\small An experiment compares $k$ treatments with $n_i$ replicates. Differences among means feed SCTR; within-treatment variation feeds SCE.\end{frame}
\begin{frame}{Application: means}\small The grand mean $\bar Y_{..}$ and treatment means $\bar Y_{i.}$ separate between-treatment from within-treatment variation.\end{frame}
\begin{frame}{Application: ratio}\small When CMTR is much larger than CME, $F$ grows and evidence against equal means increases.\end{frame}
\begin{frame}{Application: decision}\small A global rejection says that at least one mean differs; it does not identify the differing pairs.\end{frame}
\begin{frame}{Assumptions}\small Validity depends on independent errors, common variance, and approximate normality. Randomization supports independence.\end{frame}
\begin{frame}{Interpretation}\small The result is global. Report means, differences, and intervals as well as the significance decision to convey practical magnitude.\end{frame}
\begin{frame}{Communication}\small Report SCTR, SCE, degrees of freedom, mean squares, $F$, $p$-value, and the contextual conclusion.\end{frame}
\begin{frame}{Synthesis}\small\begin{itemize}\item CME estimates noise.\item CMTR combines noise and effects.\item $F$ compares them.\item The contrast is global.\end{itemize}\end{frame}
\begin{frame}{Chapter connection}\small A significant $F$ calls for multiple comparisons; the next section presents LSD, Tukey, and Bonferroni.\end{frame}
\end{document}
